\begin{textsample}{POS Dim 3 – rj – Score 20.00 – rj/rj\_inf\_cam\_exp\_6.txt}  \label{ex:f3_pos_011}
O trabalho pedagógico de cada \textbf{professora} da unidade de Educação \textbf{Infantil} ganha força e expressão na medida em que ela organiza situações e formas de estimular o desenvolvimento da autonomia \textbf{infantil} em relação a : relacionar -se com os companheiros, e conhecer -se e \textbf{cuidar} de si.
Para as \textbf{crianças} aprenderem a se relacionar com os companheiros algumas estratégias podem ser adotadas pela \textbf{professora} : - Organizar o ambiente e as \textbf{rotinas} para \textbf{acolher} as \textbf{crianças} ingressantes na unidade, ou mesmo as \textbf{crianças} já matriculadas após um período de férias, ou de adoecimento, no chamado “ processo de adaptação ”. - Estruturar um ambiente tranquilo e favorecedor do estabelecimento de interações pelas \textbf{crianças}, compreendendo a \textbf{movimentação} delas como intenções exploratórias e como forma de comunicação. - Possibilitar às \textbf{crianças} participar de ações individuais e em grupo, que as ajudem a entender os direitos e as obrigações das pessoas. - Ajudar cada \textbf{criança} a reconhecer a existência do ponto de vista do outro e considerar possíveis sentimentos, intenções e opiniões das demais pessoas, construindo atitudes negociadoras e tolerantes. - Comunicar com clareza às \textbf{crianças} instruções sobre a organização física e social do ambiente, de modo a fortalecer sua autonomia e colaboração. - Oferecer materiais e propor atividades em que as \textbf{crianças} percebam a necessidade de compartilhar e cooperar. - Ajudar as \textbf{crianças} a organizar tarefas a serem realizadas em grupo e a refletir sobre eventual quebra das regras coletivamente decididas. - Incluir as \textbf{crianças} na caracterização e arranjo dos espaços que mais frequentam e no \textbf{cuidado} de seus objetos para mantê -los bem conservados e de fácil acesso. - Atuar quando as demais \textbf{crianças} reagem a uma determinada \textbf{criança}, coibindo preconceitos, agressões, assédios, de modo a ampliar o olhar de todas para a importância de se respeitar os/as \textbf{colegas}. - \textbf{Apoiar} as \textbf{crianças} com deficiências, transtornos globais de desenvolvimento e altas habilidades/superdotação a se perceber como integrantes dos grupos \textbf{infantis}, \textbf{demonstrando} confiança em suas possibilidades de aprender com os \textbf{colegas}, estimulando -as diante de dificuldades, enquanto acompanha o que o conjunto de \textbf{crianças} pode aprender com estes \textbf{colegas}. - \textbf{Cuidar} para que os espaços, materiais, objetos e \textbf{brinquedos}, procedimentos e formas de comunicação sejam adequados às especificidades e singularidades do \textbf{brincar} e do interagir das \textbf{crianças}, em especial daquelas com deficiências, transtornos globais de desenvolvimento e altas habilidades/superdotação.
Algumas ações da \textbf{professora} podem ajudar as \textbf{crianças} a aprender a se conhecer e a \textbf{cuidar} de si.
São elas, dentre outras : - Aconchegar as \textbf{crianças} quando demandam \textbf{ajuda} ( pelo choro, pedido de \textbf{colo}, silêncio prolongado, birra ) para lidar com fortes emoções. - \textbf{Ouvir} as \textbf{crianças} e \textbf{apoiar} para que expressem seus sentimentos, planos, ideias, apontem suas \textbf{brincadeiras} e atividades preferidas e as não desejadas, e narrem suas vivências. - Incentivar as \textbf{crianças} a identificar elementos que lhes provoquem medo, apoiá -las para superá -lo e ter uma atitude ativa diante de uma dificuldade. - Ajudar as \textbf{crianças} a reconhecer em si sensações produzidas por diferentes estados fisiológicos e a comunicar que estão com sede, fome, dor, frio etc. - Comentar as ações e avaliar as produções de cada \textbf{criança} ( desenhos, esculturas, movimentos de dança, narrativas etc. ) respeitando as emoções vividas pelas \textbf{crianças}, fortalecendo sua autoestima. - Garantir igualdade no tratamento de \textbf{meninas} e \textbf{meninos}, disponibilizando todos os \textbf{brinquedos} e outros materiais para todas as \textbf{crianças} e propondo a realização de atividades em que as \textbf{crianças} possam participar independentemente do gênero. - Tratar as \textbf{crianças} e familiares pelos respectivos nomes e coibir o uso de apelidos pejorativos no tratamento dado pelas \textbf{crianças} aos \textbf{colegas} ou \textbf{adultos}. - Respeitar as diferentes formas de arranjos familiares, as opções religiosas das famílias, e \textbf{acolher} as opiniões e aspirações dos pais sobre seus filhos. - Incluir no cotidiano \textbf{brinquedos}, imagens e narrativas que promovam a construção pelas \textbf{crianças} de uma relação positiva com seus grupos de pertencimento. - Interagir com as \textbf{crianças} enquanto realizam ações de \textbf{cuidado} individual, como as trocas de fraldas, banho, sono, alimentação, de modo comunicativo e \textbf{atento}, em um ambiente planejado, seguro, aconchegante, diversificado. - \textbf{Apoiar} e incentivar as \textbf{crianças} a ter maior autonomia em relação a seu \textbf{cuidado} pessoal, tais como escovar os dentes, colocar sapatos ou o agasalho, pentear os cabelos, servir -se sozinha nas refeições e organizar seus pertences, estimulando -as a se auxiliarem nessas tarefas.

% matched lemmas: acolher, adulto, ajuda, apoiar, atento, brincadeira, brincar, brinquedo, colega, colo, criança, cuidado, cuidar, demonstrar, infantil, menino, movimentação, ouvir, professora, rotina
\end{textsample}
